% Abstract. arXiv's abstract field is plain text and limited to 1920 characters; keep this in sync
% with what you paste into the submission form (no LaTeX commands, no citations there).
\begin{abstract}
Video platforms re-encode every upload with a policy the uploader cannot observe: the codec, bit
budget and ingest filtering depend on properties of the upload itself. A common creator heuristic
exploits this by uploading at a higher resolution to obtain a better encoding tier. We ask whether
an uploader can recover quality \emph{beyond} this tier trick by optimising a pre-compensation
applied before upload, treating the platform as a non-differentiable black box. We study two
routes, sample-efficient black-box search over an interpretable filter chain and a learned
pre-compensation network trained through a differentiable emulator of the platform, and we
separate the tier gain from the filter gain. On simulated platforms, the optimised pre-compensation
improves the training objective by a large margin, but independent metrics that were not optimised
(SSIMULACRA2, XPSNR, DISTS) confirm only the tier gain; most of the apparent filter gain is colour
and contrast change that the objective rewards. Every single-metric objective we tried degrades at
least one metric it did not optimise. A learned pre-compensation exploits its emulator unless it is
constrained by a trust region. In a pre-registered design with held-out confirmation clips, a
learned pre-compensation does not beat class-tuned sharpening on natural low-motion content.
Anchoring its forward pass on the black box closes about 80\,\% of the gap, showing emulator
exploitation; the rest trades the independent metrics for its training loss. An untuned mild
sharpening applied only to natural low-motion content matches the tuned configuration, and the
upload bitrate matters more than any filter. The benefit depends on the platform's rate headroom;
validating it on a platform takes about 20 uploads.
\pending{P2 result: real-platform validation and viewer study.}
Code, clip manifests and the pre-registration are released.
\end{abstract}
